\newpage
\begin{center}
\large\textbf{Минобрнауки России}

\large\textbf{Юго-Западный государственный университет}
\vskip 1em
\normalsize{Кафедра программной инженерии}
\vskip 1em
\ifВКР{
        \begin{flushright}
        \begin{tabular}{p{.4\textwidth}}
        \centrow УТВЕРЖДАЮ: \\
        \centrow Заведующий кафедрой \\
        \hrulefill \\
        \setarstrut{\footnotesize}
        \centrow\footnotesize{(подпись, инициалы, фамилия)}\\
        \restorearstrut
        «\underline{\hspace{1cm}}»
        \underline{\hspace{3cm}}
        20\underline{\hspace{1cm}} г.\\
        \end{tabular}
        \end{flushright}
        }\fi
\end{center}
\vspace{1em}
  \begin{center}
  \large
\ifВКР{
ЗАДАНИЕ НА ВЫПУСКНУЮ КВАЛИФИКАЦИОННУЮ РАБОТУ
  ПО ПРОГРАММЕ БАКАЛАВРИАТА}
  \else
ЗАДАНИЕ НА КУРСОВУЮ РАБОТУ (ПРОЕКТ)
\fi
\normalsize
  \end{center}
\vspace{1em}
{\parindent0pt
  Студента \АвторРод, шифр\ \Шифр, группа \Группа
  
1. Тема «\Тема\ \ТемаВтораяСтрока»
\ifВКР{
утверждена приказом ректора ЮЗГУ от \ДатаПриказа\ № \НомерПриказа
}\fi.

2. Срок предоставления работы к защите \СрокПредоставления

3. Исходные данные для создания программной системы:

3.1. Перечень решаемых задач:}

\renewcommand\labelenumi{\theenumi)}

\begin{enumerate}
\item провести анализ предметной области и выявить процессы регистрации и учёта данных пользователей, подлежащие автоматизации;
\item спроектировать структуру программно-информационной системы регистрации данных пользователей в виде веб-приложения CRM;
\item разработать и реализовать подсистемы управления пользователями, клиентами, автомобилями и заявками с возможностью выполнения операций создания, просмотра, редактирования и удаления данных;
\item реализовать механизм аутентификации и авторизации пользователей с разграничением прав доступа в соответствии с ролями «Администратор», «Менеджер» и «Мастер»;
\item реализовать регламентированный жизненный цикл заявок (workflow) с проверкой допустимости переходов между состояниями и контролем прав пользователей;
\item реализовать механизм назначения мастеров на заявки и возможность принятия заявок в работу;
\item реализовать автоматическое ведение журнала изменений (аудита) по ключевым сущностям системы;
\item провести функциональное тестирование разработанной программной системы и оценить корректность её работы.
\end{enumerate}

{\parindent0pt
  3.2. Входные данные и требуемые результаты для программы:}

\begin{enumerate}
\item Входными данными для программной системы являются: данные пользователей системы (учётные записи, профили и роли); сведения о клиентах (ФИО, телефон, адрес электронной почты); данные об автомобилях (марка, модель, год выпуска, государственный регистрационный номер, принадлежность клиенту); сведения о заявках на обслуживание (описание работ, комментарии, статусы, назначенные менеджеры и мастера); данные, необходимые для аутентификации пользователей; сведения об изменениях, вносимых в ключевые сущности системы.
\item Выходными данными для программной системы являются: список и карточки пользователей системы с указанием их ролей; списки и карточки клиентов и автомобилей; перечень заявок с возможностью фильтрации и поиска; сведения о текущем состоянии и истории изменения заявок; информация о назначенных исполнителях; журнал аудита действий пользователей с возможностью фильтрации по сущности, пользователю и периоду времени; уведомления об ошибках доступа и результатах выполнения операций в интерфейсе системы.
\end{enumerate}

{\parindent0pt

  4. Содержание работы (по разделам):
  
  4.1. Введение.
  
  4.1. Анализ предметной области.
  
4.2. Техническое задание: основание для разработки, назначение разработки,
требования к программной системе, требования к оформлению документации.

4.3. Технический проект: общие сведения о программной системе, проект
данных программной системы, проектирование архитектуры программной системы, проектирование пользовательского интерфейса программной системы.

4.4. Рабочий проект: спецификация компонентов и классов программной системы, тестирование программной системы, сборка компонентов программной системы.

4.5. Заключение.

4.6. Список использованных источников.

5. Перечень графического материала:

\списокПлакатов

\vskip 2em
\begin{tabular}{p{6.8cm}C{3.8cm}C{4.8cm}}
Руководитель \ifВКР{ВКР}\else работы (проекта) \fi & \lhrulefill{\fill} & \fillcenter\Руководитель\\
\setarstrut{\footnotesize}
& \footnotesize{(подпись, дата)} & \footnotesize{(инициалы, фамилия)}\\
\restorearstrut
Задание принял к исполнению & \lhrulefill{\fill} & \fillcenter\Автор\\
\setarstrut{\footnotesize}
& \footnotesize{(подпись, дата)} & \footnotesize{(инициалы, фамилия)}\\
\restorearstrut
\end{tabular}
}

\renewcommand\labelenumi{\theenumi.}
